% Cover letter, Physical Review D.
% NOT sent. Author and affiliation intentionally unassigned.
\documentclass[11pt]{letter}
\usepackage[margin=1in]{geometry}
\usepackage{amsmath,amssymb}
\newcommand{\TODO}[1]{\textbf{[#1]}}
\signature{\TODO{author names not assigned}}
\address{\TODO{affiliation not assigned}}
\begin{document}
\begin{letter}{The Editors\\ Physical Review D}
\opening{Dear Editors,}

We submit for your consideration the manuscript \emph{``Symmetry control of
exceptional structure in the massive scalar quasinormal spectrum of doubly
rotating five-dimensional Myers--Perry black holes.''}

Exceptional points have become an active topic in black-hole spectroscopy, and
exceptional lines have recently been established in the de Sitter members of the
Myers--Perry family. It is natural to expect that the asymptotically flat,
doubly rotating case should behave similarly, and in particular that the
enhanced $U(2)$ symmetry at equal angular momenta should generate degeneracies
which unfold into exceptional points once the symmetry is broken. That is the
standard mechanism for symmetry-protected exceptional structure in non-Hermitian
systems.

Our central point is that this expectation fails, and fails for an exact
reason. Because both azimuthal Killing vectors survive for arbitrary rotation
parameters, the quasinormal pencil is a direct sum over azimuthal sectors, and
the Jordan chains of a direct sum are the union of the summands' chains. A
coincidence of frequencies from different sectors is therefore always
semisimple. The equal-spin multiplet degeneracies are protected from splitting
and forbidden from being defective by one and the same fact. We regard this
structural statement, rather than the accompanying numerics, as the main
contribution.

The remainder of the paper takes the redirection seriously. We classify the
discrete symmetry group that survives --- a Klein four-group acting on
parameters \emph{and} sector labels --- and derive from it a termination law:
in diagonal sectors an exceptional line, if one existed, would have to occur in
mirror pairs and terminate quadratically on the equal-spin surface rather than
crossing it. We then search the one channel the no-go leaves open, with a branch
atlas of $50\,558$ continued points across seven sectors, and report a bounded
exclusion with three rescaling-invariant diagnostics.

Two methodological points may interest referees independently of the physics.
First, we show that the derivative of a spectral condition cannot bound the
distance to an exceptional point, because the condition is defined only up to a
nonvanishing analytic factor; we give an invariant replacement and exhibit two
conditions with identical zeros whose derivatives differ by more than $10^5$.
Second, we report two branch-tracking defects that each produced a spurious
near-degeneracy, and the diagnostics that caught them. We include these because
we believe the failure modes are generic to this class of search.

We are explicit about what is not established. The exclusion is bounded, over a
declared domain, and is not a theorem; its minimum is attained on the boundary
of the searched box; and the quasiresonant regime is not searched, because there
the recurrence-free solvers are inapplicable by formulation rather than by
convergence failure --- a point we establish quantitatively and which we believe
corrects a common misattribution of such failures to near-extremality.

All code, machine-readable results and reproduction scripts are publicly
available, and the central result is regenerated by a single command from a
clean clone.

We suggest as possible referees \TODO{referee suggestions not assigned}.

\closing{Sincerely,}
\end{letter}
\end{document}
